\Introduction

Вопросы, связанные с питанием и профилактикой неинфекционных заболеваний, являются устойчивым направлением исследований в сфере общественного здоровья. По данным Всемирной организации здравоохранения, факторы образа жизни, включая питание, рассматриваются как значимая составляющая профилактики и снижения риска хронических заболеваний \cite{who2025stats}. В этой связи востребованы цифровые инструменты, позволяющие системно организовывать рацион и снижать долю рутинных операций при его планировании.

Целью выпускной квалификационной работы является разработка веб-приложения для планирования питания с учетом предпочтений пользователей.

Для достижения поставленной цели необходимо решить следующие задачи:
\begin{compactlist}
  \item исследовать предметную область;
  \item определить функции продукта и технологии разработки;
  \item проанализировать целевую аудиторию;
  \item проанализировать аналоги и конкурентов;
  \item спроектировать структуру базы данных;
  \item спроектировать интерфейс;
  \item разработать клиентскую часть;
  \item разработать серверную часть;
  \item разработать модули проверки макета и расчета стоимости;
  \item разместить веб-сервис на сервере.
\end{compactlist}

Объект исследования --- процесс планирования питания, включающий подбор рецептов, учет пользовательских ограничений, расчет пищевой ценности, формирование планов питания и связанных с ними списков покупок.

Предмет исследования --- методы и программные средства автоматизации планирования питания в формате веб-приложения.

Практическая значимость работы состоит в создании прикладного решения, позволяющего пользователю сократить время на планирование рациона, повысить прозрачность контроля КБЖУ и упростить формирование меню и списка покупок.

